\scanpage{139}

% L1, cancelled heading
\deleted{(Dém. plus bas)}

% L2, upper text
Reste le cas $N=2n$, \inserted{et} $H=\{n\}$ i.e. $\underline{\Delta}_{\{n\}}(E,q)$.

% L3–7, upper text
\deleted{\unsure{On}} \deleted{Soit} Supposons d'abord 2 inv. sur $S$. Soit $F\in\Delta_{\{n\}}(E,q)$, \deleted{\uline{\unsure{Proposition}}} donc $F=F^\perp$ (via lemme 1), \deleted{d'où} $E/F\xrightarrow{\sim}F^\vee$, d'où un isom. can.

% L8, display with lower brace
\[
\deleted{\omega_F\,\unsure{:}}\quad
\underbrace{\det F\otimes\det F^\vee}_{\substack{\shortparallel\\ \underline{O}_S}}\longrightarrow\det E
\]

% L9–12, middle text
qu'on interprète comme une \uline{section} de $\det E$, notée $\omega_F$. Choisissons un suppl$^{\mathrm t}$ \inserted{$G$} (sur $S$ supposé affine) de $F$ dans $E$, donc

% L13, middle display
\[
E\simeq F\oplus G,\qquad \varphi\mid F\times G,\ \text{une dualité}
\]

% L14–16, lower text
Soit $(e_1\mathbin{-}e_n)$ une base de $E$ $(\exists\ \mathrm{loc.}\ (\mathrm{Zar}))$, $e_1^*\mathbin{--}e_n^*$ la \deleted{\unsure{f}}base duale $G\simeq F^\vee$, définie par

% L17, lower display
\[
\varphi(e_i,e_j^*)\simeq\delta_{ij}
\]

% L18, lower text
on a alors

% L19, bottom display
\[
\omega_F=e_1\wedge\mathbin{--}\wedge e_n\wedge e_1^*\wedge\mathbin{-}\wedge e_n^*
\]

% bottom left, pencil
% archive sheet number: 138
